\subsection{The exact comparison with compatible elliptic pairs}

\begin{result}{E.43}{Proved}{The compatible-pair group is the graph of the elliptic involution}
Let $I\Gamma$ be the group of compatible pairs in
\cite[Definition 5.4 and Proposition--Definition 5.6]{Lochak2025Aut}.
In the conventions of this paper its elements are pairs
$A=(\lambda,f_A)$, $B=(\lambda,f_B)$ in $\widehat{\mathrm{GT}}$
whose words satisfy, for some $a,b,k\in\widehat{\mathbb Z}$,
\[
 f_B(u^2,w^2)=w^{2a}f_A(u,w^4)u^{2b}(uw^2)^{2k}
                   \quad\text{in }\widehat B_3.             \tag{E.200}
\]
Here $u,w$ are the positive standard braid generators.
Then the precise comparison is
\[
 I\Gamma=\{(A,\mathcal J(A)):A\in\widehat{\mathrm{GT}}\},
 \qquad \operatorname{pr}_0:I\Gamma
                    \xrightarrow{\sim}\widehat{\mathrm{GT}}.
                                                               \tag{E.201}
\]
The subgroup of equal compatible pairs is
$\{(A,A):\mathcal J(A)=A\}$, with its actual inclusion in
(E.201); E.44 proves that it is the whole diagonal copy
of $\widehat{\mathrm{GT}}$.  This identifies that defining pair equation; it
does not import a curve-complex automorphism theorem from
the cited paper.
\end{result}
\begin{proof}
We check the exponents and the comparison directly.  Put
$c=c_2(f_A)$ and set $z=(uw)^3$.  The braid relation gives
\[
 (uw^2)^2=uw(wuw)w=uw(uwu)w=(uw)^3=z.                       \tag{E.202}
\]
The element $z$ is central.  Pass to the central quotient
of the mixed subgroup $\langle u,w^2\rangle$ and put $v=w^2$.
It is the exact group
$R=\langle u,v\mid(uv)^2=1\rangle^{\wedge}$.
Its index-two subgroup with $u$ odd and $v$ even is free on
$v,U=u^2$.  The identities in this group give
\[
       uv^2u^{-1}=(v^{-1}U^{-1})^2.                         \tag{E.203}
\]
Indeed write $s=uv$, so $s^2=1$ and
$svs=v^{-1}U^{-1}$.  Both sides of (E.203) are $sv^2s$.

By the defining Schreier calculation of E.22, the exponent
vector of $f_A(x,y)$ in the parity subgroup with basis
$(x^2,y,xyx^{-1})$ is $(0,c,-c)$.  Under $x\mapsto u$,
$y\mapsto v^2$, its last two basis vectors have images
$(2,0)$ and $(-2,-2)$ in the abelianization with basis
$(v,U)$.  Thus
\[
       [f_A(u,v^2)]_{\langle v,U\rangle^{\rm ab}}
                         =(4c,2c).                          \tag{E.204}
\]
The left side of (E.200) is a commutator word in $U,v$.
Abelianizing (E.200) in this index-two subgroup therefore
forces $a=-4c$ and $b=-2c$.  Abelianizing the original
equation in $\widehat B_3^{\rm ab}=\widehat{\mathbb Z}$
then gives $2a+2b+6k=0$, hence $k=2c$.  Multiplication by
$6$ is injective on $\widehat{\mathbb Z}$, so no residual
central ambiguity occurs.

With those forced coefficients, (E.200) becomes
\[
 f_B(u^2,w^2)=w^{-8c}f_A(u,w^4)u^{-4c}(uw)^{6c}.             \tag{E.205}
\]
E.30 proves exactly (E.205) for $B=\mathcal J(A)$.
Conversely the two generators $u^2,w^2$ generate a free
profinite subgroup with that basis.  One direct verification
uses the pure three-strand braid presentation: writing
$x_{12}=u^2$, $x_{23}=w^2$ and
$z=x_{12}x_{13}x_{23}$, eliminate $x_{13}$.  Its defining
relations say exactly that $z$ commutes with $x_{12},x_{23}$,
so $P_3=F(x_{12},x_{23})\times\langle z\rangle$.
The free factor has full topology, and $P_3$ has finite
index in $B_3$.  The evaluation at $u^2,w^2$ is therefore
injective after full completion.  Equation (E.205) uniquely
specifies $f_B$, and yields $B=\mathcal J(A)$.

E.28 proves that $\mathcal J$ is a continuous homomorphism
and involution.  Thus $A\mapsto(A,\mathcal J(A))$ and the
first projection in (E.201) are continuous inverse group
homomorphisms.  Requiring the two entries to agree is
precisely the fixed-point equation, with the inclusion
stated after (E.201).

For the geometric convention in the cited source, its
elliptic group allows the two marked points to be exchanged.
For an ordered elliptic pair $(E,p_1,p_2)$, the map
$z\mapsto p_1+p_2-z$ gives a natural order-two isomorphism
to the exchanged pair.  Naturality shows that the resulting
label-exchange extension is the direct product of the pure
group $\Gamma_{1,2}$ with this central $C_2$.  Projection
to the pure factor and inclusion of the pure subgroup are
the actual convention maps.  The braid subgroup and words
in (E.200) lie in that pure factor; the label exchange does
not alter the displayed compatibility equation.
This is the exact comparison with Definition 5.4.  The
general curve-complex rigidity and procongruence assertions
in the later theorem of that source are separate proof
inputs and are not used here.
\end{proof}
